\documentclass[12pt]{report}

% ==============================================================================
% CONFIGURACIÓN DEL DOCUMENTO Y METADATOS (Ajusta aquí tus preferencias)
% ==============================================================================

% Selección de estilo y lengua
\newcommand{\docstyle}{dark}       % Opciones: dark, formal
\newcommand{\doclanguage}{es}      % Opciones: la, ca, es, en

% Información básica del documento, usar \textquotesingle per '
\newcommand{\titlename}{Systema Notationum}
\newcommand{\authorname}{Martí Figverola Giner}
\newcommand{\subtitle}{-}
\newcommand{\sdate}{Mayo 2026}
\newcommand{\pdate}{-}
\newcommand{\titleheader}{\titlename\ }

% Profundidad del TOC (table of contents)
\setcounter{tocdepth}{1}		% 0 = Chapter, 1 = Section, 2 = Subsection, 3 = Subsubsection

% ==============================================================================
% CARGA DE MÓDULOS DE CONFIGURACIÓN (En orden estricto de dependencias)
% ==============================================================================
\usepackage{config/01_colors}       			% Paleta de colores RGB
\usepackage{config/02_basics}       			% Márgenes, parskip, enumitem, hipervínculos
\usepackage{config/03_typography}   			% Fuentes, Polyglossia, lengua e IPA           % AQUÍ POT PETAR, CONSULTA L'ARXIU typography
\usepackage{config/04_math}         			% Paquetes base matemáticos
\usepackage{config/05_symbols}      			% Símbolos personalizados y atajos
\usepackage{config/06_tikz}         			% TikZ, CircuiTikZ y entornos fig/graph/circuit
\usepackage{config/07_toc}       				% Títulos de partes, capítulos y secciones
\usepackage{config/08_header_footnote}      	% Encabezados (fancyhdr) y reglas superiores TikZ
\usepackage{config/09_titles}    				% Pie de página, notas al pie y color global
\usepackage{config/10_tables}       			% Formato de tablas y celdas
\usepackage{config/11_tcolorbox}    			% Estilo base de cajas tcolorbox
\usepackage{config/12_theorem_environments}     % Entornos: definición, teorema, ejemplo...
\usepackage{config/13_code}         			% Bloques de código fuente
\usepackage{config/90_references}   			% Formato de referencias cruzadas
\usepackage{config/91_bibliography} 			% Fuentes y bibliografía
\usepackage{config/99_language}     			% Selección e internacionalización del idioma

% ==============================================================================
% COMANDOS ESPECÍFICOS DE ESTE DOCUMENTO
% ==============================================================================



% ==============================================================================
% CUERPO DEL DOCUMENTO
% ==============================================================================

\begin{document}
	
	\firstpage{60}	% tamaño texto
	
	\chapless{Praefātiō}
	
	Auctor operis cuj titulus est \textit{\titlename}, Mārtīnus Fīgherōla Giner, Catalānicē Martí Figuerola Giner Sicart Sobrepera Esquius Garcia Corberó Bosch [m\textturna\textfishhookr'\textsubbridge{t}i fi.\textlowering{\textgamma}\textlowering{e}'\textfishhookr\textlowering{o}.\textltilde ä \textyogh i'ne \textsubbar{s}i'kä\textfishhookr\textsubbridge{t} \textsubbar{s}\textlowering{o}.\textlowering{\textbeta}\textfishhookr \textlowering{e}'p\textlowering{e}.\textfishhookr ä \textturna \textsubbar{s}'kiw\textsubbar{s} g\textturna\textfishhookr'\textsubbar{s}i.\textturna\ k\textlowering{o}\textfishhookr.\textlowering{\textbeta}\textlowering{e}'\textfishhookr\textlowering{o} 'b\textopeno \textsubbar{s}k], Barcinōne in Hispāniā diē decimō mēnsis Augustī annō bis mīllēsimō quīntō circā hōram quīntam decimam nātus, sōlus et prīncipālis est auctor tōtīus textūs. Quīcumque error in textū invenītur, ad sōlum auctōrem pertinet ejusque est tōta rēspōnsābitās; quī, quandōcumque fierī potest, nōtificandus est.
	
	\addtoc
	
	\chap{Introducción}
	
	\section{Objetivo}
	Se presenta un sistema de metanotación para dar forma a cualquier escrito que exponga cualquier tema.
	
	\section{Lenguajes}

	
	
	\chap{Hiperlenguaje}		
	
	Primero es preciso aclarar términos básicos, luego se definirán conceptos nuevos de este texto, los hiperlenguajes y los hiperlenguajes figuerolianos.
	
	\begin{minibloque}[Definiciones]
		\begin{itemize}
			\item \textbf{Lenguaje: }Es un sistema de comunicación organizado.
			\item \textbf{Metalenguaje: }Es un lenguaje que se usa para hablar acerca de otro lenguaje.
			\item \textbf{Lenguaje Objeto: }Es un lenguaje descrito por un metalenguaje.
			\item \textbf{Lengua: }Es un lenguaje humano, pueden ser naturales o artificiales/construidas.
		\end{itemize}
	\end{minibloque}
	
	\begin{minidefin}
		Un hiperlenguaje es el metalenguaje fundamental que da estructura a un texto y describe sus conceptos.
	\end{minidefin}
	
	Para exponer cualquier tema al lector, se requiere siempre de una lengua conductora (generalmente una lengua natural, como el español) que sirva de guía. En este contexto, la lengua en la que se redacta el texto equivale al metalenguaje del propio tema, es el marco que permite formalizar y describir otras disciplinas. A esta estructura lingüística global la denominamos hiperlenguaje.
	
	La clave a destacar es que todo tema o concepto puede modelarse formalmente como un lenguaje propio, dotado de su correspondiente ortografía, sintaxis y semántica, aunque esta misma sea casi una copia del hiperlenguaje (como la filosofía o el derecho).
	
	\begin{subminibloque}[Ejemplos]
		\begin{itemize}
			\item \textbf{Matemáticas:} 
			El hiperlenguaje articula la explicación teórica y define los términos y símbolos metalingüísticamente, especialmente en lógica matemática para dar significado a expresiones de lógica de primer orden y demostrar metateoremas\footnote{Un metateorema es una proposición verdadera que se demuestra con el metalenguaje, no dentro del propio sistema formal definido matemáticamente, como el (meta-)teorema de la deducción.}.
			
			\item \textbf{Filosofía y Derecho:} 
			La ortografía y la sintaxis del tema son idénticas a las de los hiperlenguajes estándares (lenguas naturales). Sin embargo, el lenguaje del tema redefine la semántica asignando significados técnicos rigurosos a palabras comunes («ente», «jurisprudencia»...).
		\end{itemize}
	\end{subminibloque}

	
	
	
	\section{Hiperlenguajes Figuerolianos}
	
	Ahora es conveniente definir un conjunto de hiperlenguajes que se ajusten elegantemente a las necesidades de un texto, y que todos estos hiperlenguajes compartan un léxico y una sintaxis simbólica común. Para eso es necesario reemplazar y expandir los signos de puntuación de las lenguas que forman dichos hiperlenguajes.
	
	\begin{minidefin}
		Un hiperlenguaje figueroliano se compone de una lengua virgen (sin sus signos de puntuación propios) y de los hipersímbolos figuerolianos.
	\end{minidefin}
	
	Los hipersímbolos son la clave de los hiperlenguajes figuerolianos, de hecho lo único que separa cada hiperlenguaje figueroliano son las lenguas vírgenes empleadas. En otras palabras, los hipersímbolos se adaptan a cualquier lengua (con escritura estándar) para formar hiperlenguajes figuerolianos.
	
	\subsubsection{Encapsulación de Otros Lenguajes}
	Existe una aparente paradoja al analizar cómo conviven el hiperlenguaje y los lenguajes objeto dentro de un mismo enunciado. Podría pensarse que el hiperlenguaje se limita únicamente a las partículas de enlace. Sin embargo, el hiperlenguaje abarca la totalidad de la estructura del enunciado, como veremos con este ejemplo:
	\[
	\underset{\text{Hiperlenguaje}}{\underbrace{\underset{\text{Español}}{\underbrace{\vphantom{\text{Lg}}\text{La extensa}}}\text{ }\underset{\text{Filosofía}}{\underbrace{\vphantom{\text{Lg}}\text{ontología}}}\text{ }\underset{\text{Español}}{\underbrace{\vphantom{\text{Lg}}\text{puede demostrar}}}\text{ }\underset{\text{Matemáticas}}{\underbrace{\vphantom{\text{Lg}}\varnothing = \{\}}}\,\,\,\underset{\text{Hipersímbolo}}{\underbrace{\vphantom{\text{Lg}}\text{.}}}}}
	\]
	No obstante, en su forma estricta, a veces es conveniente usar el delimitador <<\textnoto{◌}>> (se verá a continuación \ref{sec:hypersymbola} REFERIR BIEN) para marcar claramente qué son otros lenguajes (no se marca <<\textnoto{◌}>> como hipersímbolo por simplicidad visual).
	\[
	\underset{\text{Hiperlenguaje}}{\underbrace{\underset{\text{Español}}{\underbrace{\vphantom{\text{Lg}}\text{La extensa}}}\text{ }\text{<<}\underset{\text{Filosofía}}{\underbrace{\vphantom{\text{Lg}}\text{ontología}}}\text{>>}\text{ }\underset{\text{Español}}{\underbrace{\vphantom{\text{Lg}}\text{puede demostrar}}}\text{ }\text{<<}\underset{\text{Matemáticas}}{\underbrace{\vphantom{\text{Lg}}\varnothing = \{\}}}\text{>>}\,\,\,\underset{\text{Hipersímbolo}}{\underbrace{\vphantom{\text{Lg}}\text{.}}}}}
	\]
	Nota también que la lengua virgen provee la matriz sintáctica base (por ejemplo, la tipología Sujeto-Verbo-Objeto), definiendo las posiciones funcionales donde se alojan los conceptos. De este modo los otros lenguajes quedan \textbf{encapsulados} bajo el hiperlenguaje como elementos sintácticos de la lengua virgen.
	\[
	\underset{\text{Hiperlenguaje}}{\underbrace{\underset{\text{Español}}{\underbrace{\vphantom{\text{Lg}}\text{La extensa}}}\text{ }\underset{\text{Español (sustantivo)}}{\underbrace{\text{<<}\underset{\text{Filosofía}}{\underbrace{\vphantom{\text{Lg}}\text{ontología}}}\text{>>}}}\text{ }\underset{\text{Español}}{\underbrace{\vphantom{\text{Lg}}\text{puede demostrar}}}\text{ }\underset{\text{Español (objeto)}}{\underbrace{\text{<<}\underset{\text{Matemáticas}}{\underbrace{\vphantom{\text{Lg}}\varnothing = \{\}}}\text{>>}}}\,\,\,\underset{\text{Hipersímbolo}}{\underbrace{\vphantom{\text{Lg}}\text{.}}}}}
	\]
	En conclusión, el hiperlenguaje no actúa como un mero conector o enlace superficial entre términos dispares. Por el contrario, opera como el metalenguaje estructural definitivo que encapsula y procesa cualquier lenguaje objeto externo dentro de la matriz sintáctica de la lengua virgen. De este modo, la totalidad del enunciado (tanto la arquitectura lógica como las partículas importadas) queda plenamente integrada bajo la jerarquía del hiperlenguaje.


	
	\section{Hipersímbolos Figuerolianos}\label{sec:hypersymbola}
	
	\begin{minidefin}
		Los hipersímbolos figuerolianos se componen de:
		\begin{itemize}
			\item Hiperconstantes
			\item Hiperbloques
			\item Hipersignos
			\item Hiperdelimitadores
		\end{itemize}
	\end{minidefin}
	
	\subsection{Generalización de los Símbolos Tradicionales}
	
	Estos son los hipersímbolos tradicionales con sus respectivos usos normalemente, algunos se les extiende el uso. ASI SE PERMITE UNA CAPA PRACTICA Y OTRA CANONICA, la practica para uso mas barriobajero, y cuando se dispone del inventario simbolico precision canonica.
	\begin{tab}{3}{|+c|-C|-C|}{Símbolo & Significado & Significado Expandido}
		. & Final de frase & Final \\\hline
		, & Separación de elementos & Separador \\\hline
		; & Separación de frases & Separador superior \\\hline
		: & Introducción & Introducción \\\hline
		... & Elipsis & Elipsis \\\hline
		¿\textnoto{◌}? & Interrogación & Interrogación \\\hline
		¡\textnoto{◌}! & Exclamación & Exclamación \\\hline
		(\textnoto{◌}) & Aclaración/Comentario & Aclaración/Comentario \\\hline
		[\textnoto{◌}] & Adición a una cita, referencia, fonética... & Adición pura \\\hline
		\guillemotleft\textnoto{◌}\guillemotright & Delimitador primario & Delimitador de lenguaje \\\hline
		\textquotedblleft\textnoto{◌}\textquotedblright & Delimitador secundario & Delimitador literal \\\hline
		\textquoteleft\textnoto{◌}\textquoteright & Delimitador terciario & Delimitador parafraseador \\\hline
	\end{tab}
	
	
	Abreviaciones con punto hacerlas como: $\overline{\text{ıe}}$, $\overline{\text{eg}}$. Podriamos decir que overline es un signo que hace reune abreviaciones.
	
	YO HARIA QUE, el unico hipersímbolo sea << y >> con un argumento opcional tipo <<$_\mathcal{A}$Hola que tal>> o <<$^\mathcal{A}$Hola que tal>>.

	\begin{tab}{2}{|+c|-C|}{Signo & Significado}
		\textnoto{◌} & Marca un argumento \\\hline
		$\cdots$ & Marca continuación o elipsis \\\hline
	\end{tab}
	
	
	\guilsinglleft\textnoto{◌}\guilsinglright, \guillemotleft\textnoto{◌}\guillemotright, \textquotedblleft\textnoto{◌}\textquotedblright, \textquoteleft\textnoto{◌}\textquoteright, (\textnoto{◌}), $\llparenthesis$\textnoto{◌}$\rrparenthesis$, [\textnoto{◌}], $[$\textnoto{◌}$]$, $\llbracket$\textnoto{◌}$\rrbracket$, $\lfloor$\textnoto{◌}$\rfloor$, $\lceil$\textnoto{◌}$\rceil$, $\llfloor$\textnoto{◌}$\rrfloor$, $\llceil$\textnoto{◌}$\rrceil$, \textlangle\textnoto{◌}\textrangle, $\ulcorner$\textnoto{◌}$\urcorner$, $\llcorner$\textnoto{◌}$\lrcorner$, \textlquill\textnoto{◌}\textrquill
	
	-, Relación; --, Interrupción.
	
	\begin{tab}{2}{|+c|-C|}{Signo & Significado}
		\guilsinglleft\textnoto{◌}\guilsinglright & Delimita un metalenguaje \\\hline
		\guillemotleft\textnoto{◌}\guillemotright & Delimita un lenguaje \\\hline
		\textquotedblleft\textnoto{◌}\textquotedblright & Cita literal \\\hline
		\textquoteleft\textnoto{◌}\textquoteright & Parafraseo o cita traducida \\\hline
		(\textnoto{◌}) & Comentario \\\hline
		$[$\textnoto{◌}$]$ & Adiciones (fonética o desambiguar citas) \\\hline
		$\llbracket$\textnoto{◌}$\rrbracket$ & Referenciar \\\hline
		$\lfloor$\textnoto{◌}$\rfloor$ & Adiciones a citas literales \\\hline
		\textlangle\textnoto{◌}\textrangle & Delimita algo opcional \\\hline
		$\ulcorner$\textnoto{◌}$\urcorner$ & Nombre \\\hline
		\{\textnoto{◌}\} & Agrupación \\\hline
	\end{tab}
	
	Al analizar el término \guillemotleft Wahrheit\guillemotright\ $\llparenthesis$ˈvaːɐ̯haɪ̯t$\rrparenthesis$ $\lfloor$Frege, 1918, p. 58$\rfloor$, el \guilsinglleft$\wedge$\guilsinglright\ toma la cita literal \textquotedblleft Der Gedanke\textquotedblright\ [redactada en su borrador final] y la asimila como \textquoteleft el contenido objetivo del pensamiento\textquoteright\ (principio central de su ontología). A través de este proceso, el término formalizado $\ulcorner\nu\urcorner$ evalúa el lenguaje objeto \guillemotleft$\forall x (P(x) \to Q(x))$\guillemotright\ para computar su denotación exacta $\llbracket \text{Wahrheit} \rrbracket = 1$.
	
	
	
	\chap{Lenguaje Matemático}
	\section{Fórmulas}
	
	Las fórmulas proposicionales se componen por 3 elementos distintos:
	\begin{e1}
		\item Variables: $p_1, p_2, p_3...$
		\item Connectores: $¬, \wedge, \vee, \oplus, \to, \leftrightarrow$
		\item Paréntesis: $(, )$
	\end{e1}
	
	\begin{defin}
		Una \textbf{fórmula proposicional} se define inductivamente:
		\begin{ne1}
			\item $\forall i\in\mathbb{Z}_{\ge 0}, p_i$.
			\item Siendo $A$ una fórmula proposicional, $¬A$ también lo es.
			\item Siendo $A$ i $B$ fórmulas proposicionales, $A\otimes B$ también lo es para cualquier connector $\otimes$.
		\end{ne1}
	\end{defin}
	
	\begin{defin}
		Una fórmula $\varphi$ es \textbf{valida} o \textbf{tautológica} si $\vDash \phi$.
	\end{defin}
	\begin{defin}
		Una fórmula $\varphi$ es \textbf{demostrable} si $\vdash \phi$.
	\end{defin}
	
	
	
	
	
	
	\chap{Símbolos Figuerolianos}\label{chap:symbola}
	
	\section{Sagittulae Fīgherōliānae}
	
	

	
	
	\section{Signa Fīgherōliāna}
	
	
	
	
	
	
	
	
	\chapless{Glossary}
	
	\section*{List of Symbols}
	%	\begin{xltabular}{\textwidth}{lcX}
		%		$x$ \glsarrow Position (\siless{\metre})\\
		%		$t$ \glsarrow Time (\siless{\second})\\
		%		$v$ \glsarrow Velocity (\siless{\metre/\second})\\
		%		$a$ \glsarrow Acceleration (\siless{\metre/\second^2})\\
		%		$m$ \glsarrow Mass (\siless{\kilo\gram})\\
		%		$E$ \glsarrow Energy (\siless{\joule})\\
		%		$F$ \glsarrow Force (\siless{\newton})\\
		%		$v_s$ \glsarrow Velocity of sound (343\si{\metre/\second})\\
		%		$v_l$ \glsarrow Velocity of light ($3\cdot10^8$\si{\metre/\second})
		%	\end{xltabular}
	
	\section*{List of Figures}
	
	
	\section*{List of Tables}
	
	
	\section*{List of Codes}
	
	\section*{Acronyms}
	\begin{xltabular}{\textwidth}{lcX}
		IPA \glsarrow International Phonetic Alphabet\\
		FIPA \glsarrow Figuerola International Phonetic Alphabet\\
		FIPAS \glsarrow Figuerola International Phonetic Alphabet Simplified\\
		FIPAM \glsarrow Figuerola International Phonetic Alphabet Modified\\
		FIPAX \glsarrow Figuerola International Phonetic Alphabet Extra
	\end{xltabular}
	
	\section*{Terms}
	\renewcommand{\arraystretch}{1.5}
	\begin{xltabular}{\textwidth}{lcX}
		IPA \glsarrow International Phonetic Alphabet\\
		FIPA \glsarrow Figuerola International Phonetic Alphabet a lot of texxxxxxxxxxt for youuuuuuu sdjflasdf odsjf lasd jd \\
		FIPAS \glsarrow Figuerola International Phonetic Alphabet Simplified\\
		FIPAM \glsarrow Figuerola International Phonetic Alphabet Modified\\
		IPA \glsarrow International Phonetic Alphabet\\
		FIPA \glsarrow Figuerola International Phonetic Alphabet a lot of texxxxxxxxxxt for youuuuuuu sdjflasdf odsjf lasd jd \\
		FIPAS \glsarrow Figuerola International Phonetic Alphabet Simplified\\
		FIPAM \glsarrow Figuerola International Phonetic Alphabet Modified\\
		IPA \glsarrow International Phonetic Alphabet\\
		FIPA \glsarrow Figuerola International Phonetic Alphabet a lot of texxxxxxxxxxt for youuuuuuu sdjflasdf odsjf lasd jd \\
		FIPAS \glsarrow Figuerola International Phonetic Alphabet Simplified\\
		FIPAM \glsarrow Figuerola International Phonetic Alphabet Modified\\
		FIPAX \glsarrow Figuerola International Phonetic Alphabet Extra
	\end{xltabular}
	\renewcommand{\arraystretch}{1}
	
	\renewcommand\bibname{References}
	\bibliographystyle{unsrtnat}
	\bibliography{References}
	\addcontentsline{toc}{chapter}{References}
	
	% google scholar > cite > BibTeX
	
	
	\lastpage{1} 	% 0 --> no comercial, 1 --> sí comercial

\end{document}